\chapter{Validation et résultats}

\section{Stratégie de test}

La validation distingue les fonctions unitaires, le retrieval, les règles de sécurité et l'intégration Neo4j. Les tests couvrent notamment :

\begin{itemize}
  \item la résolution de classes et le choix du mode ;
  \item la vérification des citations ;
  \item les corrections directionnelles bornées ;
  \item le rejet des auto-relations ;
  \item la conservation de plusieurs preuves ;
  \item le refus d'importer une proposition non validée ;
  \item la récupération des chunks de preuve métier sans embedding ;
  \item l'isolation des conversations.
\end{itemize}

La dernière exécution effectuée pendant la rédaction a obtenu 49 tests réussis. Un avertissement de dépréciation Starlette reste à traiter lors d'une future mise à jour des dépendances.

\capture{tests-49.png}{Résultat de la suite de 49 tests}{fig:tests}

\section{Évaluation fonctionnelle du GraphRAG}

Le module \code{arianes/evaluation.py} regroupe des scénarios couvrant les questions structurelles, documentaires, hybrides et hors domaine. L'évaluation du retrieval reste indépendante de la disponibilité ponctuelle du fournisseur de génération.

\section{Pilote métier}

\begin{table}[H]
\centering
\caption{État actuel de la couche métier dans Neo4j}
\begin{tabularx}{0.8\textwidth}{Y>{\raggedleft\arraybackslash}p{2cm}}
\toprule
Élément & Nombre \\
\midrule
Concepts métier & 255 \\
Relations métier actives & 347 \\
Relations VALIDE & 123 \\
Relations PROBABLE & 224 \\
Relations REJETE importées & 0 \\
\bottomrule
\end{tabularx}
\end{table}

L'import repose sur des identifiants canoniques et sur \code{MERGE}. Les relations conservent leurs preuves et leurs statuts. Le retriever distingue ces relations métier des associations UML et peut calculer des chemins séparément dans les deux couches.

\section{Discussion}

Les résultats démontrent la faisabilité de la chaîne complète, depuis la source documentaire jusqu'à la réponse sourcée. Ils ne constituent pas une validation ferroviaire de l'ensemble des relations extraites. La précision d'une relation dépend de la qualité du passage, du schéma de prédicats, du prompt et de la relecture.

\section{Limites}

\begin{itemize}
  \item la sélection initiale des classes candidates du pipeline métier est principalement lexicale ;
  \item les noms CamelCase et les synonymes peuvent réduire le rappel ;
  \item les relations \code{PROBABLE} exigent une revue experte ;
  \item le pilote global documentaire n'est pas une certification du corpus ;
  \item les réponses dépendent de la disponibilité du fournisseur choisi ;
  \item les métriques techniques ne remplacent pas une évaluation par des experts du domaine.
\end{itemize}

\section{Prototype de reranking}

Un prototype de reranking supervisé a été réalisé pour étudier une piste de machine learning. Avec seulement 40 exemples déséquilibrés, il a obtenu des résultats inférieurs à la fusion hybride : nDCG@5 de 0,8725 contre 0,9369, et Recall@5 de 0,6367 contre 0,7033. Il reste donc désactivé. Cette expérience indique qu'un jeu de questions--chunks annoté par des experts est nécessaire avant toute activation.
